% Data-grounded draft
\section{Deep Learning and Ablation Results}
\label{sec:results_dl}

\subsection{Tier-wise performance}
Table~\ref{tab:dl_all} lists held-out split results for all ten Phase~3 models.
Within Tier~0, VGG16 achieves the highest macro-F1 ($0.950$) and accuracy ($0.956$), followed by EfficientNet-B0 (F1 $=0.940$) and ResNet50 (F1 $=0.905$).
ViT-B/16 lags at F1 $=0.871$, consistent with transformer data hunger on 1{,}500 images.

Tier~1 hybrids cluster between F1 $=0.910$--$0.917$: WaveletCNN (DWT branch) leads at $0.917$, followed by CNN+HOG ($0.912$), LBPNet ($0.912$), and CNN+GLCM ($0.911$).
Tier~2 fusion models underperform every Tier~1 single-branch variant: MultiBranchFusion F1 $=0.895$ and AttentionFusion F1 $=0.886$.
This ordering contradicts the naive assumption that more branches always help; instead, fusion introduces optimisation difficulty or redundant/conflicting gradients when all four descriptors are stacked without selection.

% Auto-generated by tables_to_latex.py — do not edit by hand
\begin{table}[!t]
\centering
\caption{Phase 3 deep learning models on the held-out split.}
\label{tab:dl_all}
\begin{tabular}{@{}llccc@{}}
\toprule
Tier & Model & Acc & F1-macro & Time (s) \\
\midrule
Tier 0 & EfficientNet-B0 & 0.9463 & 0.9400 & 395.9 \\
Tier 0 & ResNet50 & 0.9137 & 0.9055 & 120.4 \\
Tier 0 & VGG16 & 0.9560 & 0.9502 & 299.6 \\
Tier 0 & ViT-B/16 & 0.8827 & 0.8708 & 271.0 \\
Tier 1 & CNN+GLCM & 0.9186 & 0.9107 & 0.0 \\
Tier 1 & CNN+HOG & 0.9218 & 0.9123 & 248.9 \\
Tier 1 & LBPNet & 0.9169 & 0.9121 & 123.5 \\
Tier 1 & WaveletCNN & 0.9235 & 0.9172 & 151.5 \\
Tier 2 & AttentionFusion & 0.8974 & 0.8856 & 331.4 \\
Tier 2 & MultiBranchFusion & 0.9039 & 0.8952 & 288.1 \\
\bottomrule
\end{tabular}
\end{table}


\subsection{Feature-branch ablation}
Table~\ref{tab:ablation} isolates branch contributions relative to the best Tier~1 hybrid (WaveletCNN, $\Delta$F1 $=0$).
Removing all handcrafted branches and retaining only the VGG16 Tier~0 baseline yields F1 $=0.950$, showing that transfer learning remains strong on this dataset.
Among hybrids, attaching DWT alone matches the best single-branch score, whereas GLCM, LBP, and HOG branches incur $\Delta$F1 of $-0.007$ to $-0.005$.
Aggregating all four branches (MultiBranchFusion) degrades F1 by $0.022$ relative to WaveletCNN---the largest penalty in the ablation.

% Auto-generated by tables_to_latex.py — do not edit by hand
\begin{table}[!t]
\centering
\caption{Feature-branch ablation (held-out split). CNN-only baseline uses Tier-0 VGG16 until dedicated run completes.}
\label{tab:ablation}
\begin{tabular}{@{}lcccccc@{}}
\toprule
Model & GLCM & LBP & DWT & HOG & F1 & $\Delta$F1 \\
\midrule
CNN only (VGG16) & --- & --- & --- & --- & 0.9502 & --- \\
CNN+GLCM & \checkmark & --- & --- & --- & 0.9107 & -0.0065 \\
CNN+LBP & --- & \checkmark & --- & --- & 0.9121 & -0.0051 \\
CNN+DWT & --- & --- & \checkmark & --- & \textbf{0.9172} & 0.0000 \\
CNN+HOG & --- & --- & --- & \checkmark & 0.9123 & -0.0049 \\
All branches & \checkmark & \checkmark & \checkmark & \checkmark & 0.8952 & -0.0220 \\
\bottomrule
\end{tabular}
\end{table}


\begin{figure}[!t]
\centering
\includegraphics[width=\linewidth]{phase3_tier_comparison}
\caption{Macro-F1 by model on the held-out split, coloured by tier.}
\label{fig:tier_bar}
\end{figure}

Training-time differences are substantial: ResNet50 converges in ${\sim}120$\,s versus ${\sim}5{,}078$\,s for ViT-B/16 on the same hardware, informing deployment trade-offs independent of accuracy.
